%% kapa: exam=1 round=21 year=2010 period=2 form=A subject=appraisal_law subject_ko=부동산관계법규 seq=7 number=47
%% source: Q-Net 감정평가사 기출문제 (cst00309) artlSeq=5024181 file=extracted/1008/5024181/2010년 제21회 감정평가사 1차시험 기출문제/2교시.hwp
%% conversion: kordoc 4.24.51 -> tools/gen1.py (auto); review=proofread-gpt-sol; flags=none
\begin{question}{47}
국토의 계획 및 이용에 관한 법령상 시가화조정구역에 관한 설명으로 옳지 않은 것은?
\begin{choices}
  \item 계획관리지역에는 지정될 수 없다.
  \item 지정권자는 국토해양부장관이다.
  \item 시가화 유보기간은 도시관리계획으로 정한다.
  \item 시가화 유보기간은 5년 이상 20년 이내이다.
  \item 구역지정의 실효고시는 관보에 게재하는 방법에 의한다.
\end{choices}
\end{question}
\answer{1}
